\documentclass[11pt]{article}
\usepackage[utf8]{inputenc}
\usepackage[T1]{fontenc}
\usepackage{lmodern}
\usepackage{amsmath,amssymb}
\usepackage[letterpaper,margin=2.3cm]{geometry}
\setlength{\parindent}{0pt}
\pagestyle{plain}
\begin{document}
\textbf{Sistemas Dinámicos II}\\
\textbf{Tercer examen parcial -- Simulacro 1}

\vspace{0.6cm}
Nombre: \underline{\hspace{9cm}} \quad Clave: \underline{\hspace{3cm}}

\vspace{0.5cm}
\begin{center}
JUSTIFICA CON DETALLE TUS RESPUESTAS\\
NO se permiten libros, apuntes, calculadoras, celulares o tabletas\\
Tiempo: 90 minutos
\end{center}
\vspace{-0.2cm}
\rule{\textwidth}{0.4pt}
\vspace{0.2cm}

\begin{enumerate}
\item (2.5 puntos) Dibujar el retrato fase del siguiente sistema de ecuaciones diferenciales. Analizar la estabilidad de todos los puntos de equilibrio. Encontrar el $\omega$-límite y el $\alpha$-límite para cualquier punto en el plano.
\begin{eqnarray*}
\dot x &=& -y+x(1-x^2-y^2)(4-x^2-y^2), \\
\dot y &=& x+y(1-x^2-y^2)(4-x^2-y^2).
\end{eqnarray*}
\vspace{0.15cm}

\item (2.5 puntos) Usando el método del potencial, determinar la energía potencial $U$ del siguiente sistema así como los puntos de equilibrio, determinar la estabilidad de estos puntos de equilibrio y dibujar el retrato fase correspondiente.
$$\ddot x=-\sin 2x-\sin x.$$
\vspace{0.15cm}

\item (2.5 puntos) Encontrar los valores de bifurcación del siguiente sistema de EDO, donde $\mu\in\mathbb{R}$. Dibujar los retratos fase correspondientes y el diagrama de bifurcación (radio de las órbitas periódicas contra $\mu$).
\begin{eqnarray*}
\dot x &=& \mu x-y+x(x^2+y^2)-x(x^2+y^2)^2, \\
\dot y &=& x+\mu y+y(x^2+y^2)-y(x^2+y^2)^2.
\end{eqnarray*}
\vspace{0.15cm}

\item (2.5 puntos) Considera el sistema $\dot x=f(x)$, $x\in\Omega\subset\mathbb{R}^n$, con $\Omega$ abierto y $f\in C^1(\Omega)$. Sea $\varphi_t$ su flujo y $p\in\Omega$ tal que $\varphi_t(p)$ está definido para todo $t\ge0$ y $\lim_{t\to\infty}\varphi_t(p)=q\in\Omega$. Demostrar que $q$ es un punto de equilibrio y que $\omega(p)=\{q\}$.
\vspace{0.15cm}

\end{enumerate}

\vspace{0.5cm}
\begin{center}
\Large Buena suerte!!
\end{center}

\newpage
\textbf{Sistemas Dinámicos II}\\
\textbf{Tercer examen parcial -- Simulacro 2}

\vspace{0.6cm}
Nombre: \underline{\hspace{9cm}} \quad Clave: \underline{\hspace{3cm}}

\vspace{0.5cm}
\begin{center}
JUSTIFICA CON DETALLE TUS RESPUESTAS\\
NO se permiten libros, apuntes, calculadoras, celulares o tabletas\\
Tiempo: 90 minutos
\end{center}
\vspace{-0.2cm}
\rule{\textwidth}{0.4pt}
\vspace{0.2cm}

\begin{enumerate}
\item (2.5 puntos) Demostrar que la ecuación diferencial de segundo orden
$$\ddot z+(4z^2+\dot z^2-4)\dot z+z=0$$
tiene una solución periódica no trivial. Encontrar un anillo $\{r_1\le |(z,\dot z)|\le r_2\}$ que la contenga.
\vspace{0.15cm}

\item (2.5 puntos) (a) Encontrar la solución maximal del P.V.I. $\dot x=1+x^2$, $x(0)=0$ y su intervalo maximal de definición. Explicar por qué esto no contradice el teorema de existencia y unicidad, si $f(x)=1+x^2$ es de clase $C^\infty$ en todo $\mathbb{R}$.

(b) Para el P.V.I. $\dot x=x^2$, $x(0)=x_0$, determinar el intervalo maximal de definición de la solución en función de $x_0$.
\vspace{0.15cm}

\item (2.5 puntos) Considera el siguiente sistema, donde $\lambda>0$. Determinar las ceroclinas, encontrar los puntos de equilibrio en el primer cuadrante y su estabilidad en función de $\lambda$. Encontrar los valores de bifurcación y esbozar los retratos fase correspondientes.
\begin{eqnarray*}
\dot x &=& x(2-x-y), \\
\dot y &=& y(\lambda-x-2y).
\end{eqnarray*}
\vspace{0.15cm}

\item (2.5 puntos) Sea $V:\mathbb{R}^n\to\mathbb{R}$ de clase $C^2$ y considera el sistema gradiente $\dot x=-\nabla V(x)$.

(a) Demostrar que el sistema no tiene órbitas periódicas no triviales.

(b) Sea $p$ tal que la semiórbita positiva de $p$ es acotada. Demostrar que $\omega(p)$ está formado únicamente por puntos de equilibrio.
\vspace{0.15cm}

\end{enumerate}

\vspace{0.5cm}
\begin{center}
\Large Buena suerte!!
\end{center}

\newpage
\textbf{Sistemas Dinámicos II}\\
\textbf{Tercer examen parcial -- Simulacro 3}

\vspace{0.6cm}
Nombre: \underline{\hspace{9cm}} \quad Clave: \underline{\hspace{3cm}}

\vspace{0.5cm}
\begin{center}
JUSTIFICA CON DETALLE TUS RESPUESTAS\\
NO se permiten libros, apuntes, calculadoras, celulares o tabletas\\
Tiempo: 90 minutos
\end{center}
\vspace{-0.2cm}
\rule{\textwidth}{0.4pt}
\vspace{0.2cm}

\begin{enumerate}
\item (2.5 puntos) Encontrar el valor de bifurcación $\lambda_0$ de la siguiente EDO. Dibujar los respectivos retratos fase para $\lambda<\lambda_0$, $\lambda=\lambda_0$ y $\lambda>\lambda_0$. Así como el diagrama de bifurcación.
$$\ddot x=\lambda x-x^3.$$
\vspace{0.15cm}

\item (2.5 puntos) Usando el método del potencial, determinar la energía potencial $U$ del siguiente sistema así como los puntos de equilibrio, determinar la estabilidad de estos puntos de equilibrio y dibujar el retrato fase correspondiente.
$$\ddot x=-x^3+4x^2-3x.$$
\vspace{0.15cm}

\item (2.5 puntos) Determinar si los siguientes sistemas tienen o no órbitas periódicas.

(a) En el primer cuadrante $x>0$, $y>0$ (sugerencia: $B(x,y)=1/(xy)$),
\begin{eqnarray*}
\dot x &=& x(3-x-2y), \\
\dot y &=& y(2-x-y).
\end{eqnarray*}
(b) Demostrar que las soluciones de $\ddot z+4z-z^3=0$ con condiciones iniciales cercanas al origen son periódicas. ¿Para qué valores de la energía lo son?
\vspace{0.15cm}

\item (2.5 puntos) Considera el sistema $\dot x=f(x)$, $x\in\Omega\subset\mathbb{R}^n$, con $\Omega$ abierto y $f\in C^1(\Omega)$. Sea $K\subset\Omega$ compacto no vacío y positivamente invariante. Sea $p\in K$. Demostrar que $\omega(p)$ es no vacío y compacto.
\vspace{0.15cm}

\end{enumerate}

\vspace{0.5cm}
\begin{center}
\Large Buena suerte!!
\end{center}

\newpage
\textbf{Sistemas Dinámicos II}\\
\textbf{Tercer examen parcial -- Simulacro 4}

\vspace{0.6cm}
Nombre: \underline{\hspace{9cm}} \quad Clave: \underline{\hspace{3cm}}

\vspace{0.5cm}
\begin{center}
JUSTIFICA CON DETALLE TUS RESPUESTAS\\
NO se permiten libros, apuntes, calculadoras, celulares o tabletas\\
Tiempo: 90 minutos
\end{center}
\vspace{-0.2cm}
\rule{\textwidth}{0.4pt}
\vspace{0.2cm}

\begin{enumerate}
\item (2.5 puntos) Usando el método del potencial, determinar la energía potencial $U$ del siguiente sistema así como los puntos de equilibrio, determinar la estabilidad de estos puntos de equilibrio y dibujar el retrato fase correspondiente.
$$\ddot x=\cos x\,(2\sin x-1).$$
\vspace{0.15cm}

\item (2.5 puntos) Considera el sistema
\begin{eqnarray*}
\dot x &=& -x^3-3y, \\
\dot y &=& x-y^3.
\end{eqnarray*}
(a) Mostrar que la linealización no permite decidir la estabilidad del origen.

(b) Encontrar $a,b>0$ tales que $V(x,y)=ax^2+by^2$ sea una función de Lyapunov estricta para $(0,0)$ y concluir sobre su estabilidad.
\vspace{0.15cm}

\item (2.5 puntos) El siguiente sistema modela una presa $x$ y un depredador $y$, con $\lambda>0$. Encontrar los puntos de equilibrio con $x,y\ge0$ y su estabilidad. Encontrar el valor de bifurcación $\lambda_0$ y esbozar los retratos fase para $\lambda<\lambda_0$, $\lambda=\lambda_0$ y $\lambda>\lambda_0$. ¿Para qué valores de $\lambda$ el equilibrio donde coexisten las dos especies es un foco?
\begin{eqnarray*}
\dot x &=& x(2-x)-xy, \\
\dot y &=& y(x-\lambda).
\end{eqnarray*}
\vspace{0.15cm}

\item (2.5 puntos) (a) Sea $D\subset\mathbb{R}^2$ una región simplemente conexa y $f=(P,Q)$ un campo de clase $C^1$ en $D$ con $\partial P/\partial x+\partial Q/\partial y>0$ en todo $D$. Demostrar que el sistema $\dot x=f(x)$ no tiene órbitas periódicas contenidas en $D$.

(b) Usar (a) para demostrar que el siguiente sistema no tiene órbitas periódicas.
\begin{eqnarray*}
\dot x &=& x+y^2, \\
\dot y &=& x^2+y^3.
\end{eqnarray*}
\vspace{0.15cm}

\end{enumerate}

\vspace{0.5cm}
\begin{center}
\Large Buena suerte!!
\end{center}

\newpage
\textbf{Sistemas Dinámicos II}\\
\textbf{Tercer examen parcial -- Simulacro 5}

\vspace{0.6cm}
Nombre: \underline{\hspace{9cm}} \quad Clave: \underline{\hspace{3cm}}

\vspace{0.5cm}
\begin{center}
JUSTIFICA CON DETALLE TUS RESPUESTAS\\
NO se permiten libros, apuntes, calculadoras, celulares o tabletas\\
Tiempo: 90 minutos
\end{center}
\vspace{-0.2cm}
\rule{\textwidth}{0.4pt}
\vspace{0.2cm}

\begin{enumerate}
\item (2.5 puntos) Considera el siguiente sistema de EDO escrito en coordenadas polares, con $a\in\mathbb{R}$. Encontrar los valores de bifurcación. Dibujar todos los posibles retratos fase. Tomar $\theta\in[0,2\pi]$.
\begin{eqnarray*}
\dot r &=& r(1-r^2), \\
\dot\theta &=& a-\sin\theta.
\end{eqnarray*}
\vspace{0.15cm}

\item (2.5 puntos) Encontrar el valor de bifurcación $\lambda_0$ de la siguiente EDO. Dibujar los respectivos retratos fase para $\lambda<\lambda_0$, $\lambda=\lambda_0$ y $\lambda>\lambda_0$. Así como el diagrama de bifurcación.
$$\ddot x=x^2-2x+\lambda.$$
\vspace{0.15cm}

\item (2.5 puntos) Dibujar el retrato fase del siguiente sistema de ecuaciones diferenciales. Analizar la estabilidad de todos los puntos de equilibrio. Encontrar el $\omega$-límite y el $\alpha$-límite para cualquier punto en el plano.
\begin{eqnarray*}
\dot x &=& 4x+3y-x(x^2+y^2), \\
\dot y &=& -3x+4y-y(x^2+y^2).
\end{eqnarray*}
\vspace{0.15cm}

\item (2.5 puntos) Considera el siguiente PVI
$$(t^2-2t-8)\,\frac{d^3x}{dt^3}+e^t\,\frac{d^2x}{dt^2}+(\cos t)\,x=0,\qquad x(1)=0,\ x'(1)=2,\ x''(1)=-1.$$
(a) Demostrar que la solución existe y es única. Determinar el intervalo maximal donde el teorema garantiza que esta solución está definida.

(b) ¿Cuál sería ese intervalo si las mismas condiciones iniciales se dieran en $t_0=5$?
\vspace{0.15cm}

\end{enumerate}

\vspace{0.5cm}
\begin{center}
\Large Buena suerte!!
\end{center}

\end{document}
